\section{课程总览：自我改进是一个系统闭环}

课程从语言模型扩展到 agent：agent 有目标、环境、工具、状态和反馈，并通过多步规划与自我修正完成任务。前八讲反复出现同一个闭环：

$$
\text{Propose}\rightarrow \text{Act/Search}\rightarrow
\text{Verify}\rightarrow \text{Learn or Revise}.
$$

Test-time scaling 在参数不变时增加采样、搜索和验证；train-time scaling 把成功轨迹与 reward 写回参数；open-ended search 则进一步允许系统探索预先没有列出的解法或任务。

\videofigure{figures/lecture-roadmap.jpg}{最终讲聚焦四个未来方向：多 agent reasoning diversity、数学 meta-verification、零外部数据的任务自生成，以及 intelligence per watt。}{00:03:58--00:05:08}

\videofigure{figures/improvement-bottlenecks.jpg}{下一代自我改进要突破 reasoning 多样性、robust verification 和训练任务静态选择三类瓶颈。}{00:05:08--00:07:28}

\begin{importantbox}{自我改进的三个稀缺资源}
第一是新颖且正确的 reasoning trajectory，第二是可信 verifier，第三是处在当前能力边界上的训练任务。缺少任何一个，闭环都会早早平台化或被 reward hacking 破坏。
\end{importantbox}

\begin{knowledgebox}{Agent 是 LLM 与外部机制的联合系统}
当前 LLM 往往不足以单独承担完整目标。实际系统会编排多个模型、工具、搜索器、judge 和 memory；能力来自组件与反馈路径的共同设计。
\end{knowledgebox}

\subsection{本章小结}

课程的长期价值不在某篇最新论文，而在一组可复用原语：生成、搜索、工具、验证、reward、记忆、回溯和训练。最终讲用四项研究展示这些原语将如何继续演化。

